\documentclass[twocolumn,11pt]{article}

% Packages
\usepackage[utf8]{inputenc}
\usepackage[T1]{fontenc}
\usepackage{textcomp}
\usepackage[margin=1in]{geometry}
\usepackage{graphicx}
\usepackage{amsmath}
\usepackage{amssymb}
\usepackage{hyperref}
\usepackage{xcolor}
\usepackage{booktabs}
\usepackage{caption}
\usepackage{float}
\usepackage{parskip}
\usepackage{algorithm}
\usepackage{algorithmic}
\usepackage{titlesec}
\usepackage[numbers,sort&compress]{natbib}

% Paragraph spacing: 6pt after with stretchability
\setlength{\parskip}{6pt plus 2pt minus 1pt}
\setlength{\parindent}{0pt}

% Section spacing: 6pt before and after with stretchability
\titlespacing*{\section}{0pt}{6pt plus 2pt minus 1pt}{6pt plus 2pt minus 1pt}
\titlespacing*{\subsection}{0pt}{6pt plus 2pt minus 1pt}{6pt plus 2pt minus 1pt}
\titlespacing*{\subsubsection}{0pt}{6pt plus 2pt minus 1pt}{6pt plus 2pt minus 1pt}

% Figure and table spacing: 6pt before and after with stretchability
\setlength{\floatsep}{6pt plus 2pt minus 1pt}
\setlength{\textfloatsep}{6pt plus 2pt minus 1pt}
\setlength{\intextsep}{6pt plus 2pt minus 1pt}

% Make all pages fill to the bottom
\flushbottom
\setlength{\topskip}{1.6\topskip}

% Add more stretchability for bibliography
\setlength{\bibsep}{0pt plus 3pt minus 0pt}

% Title and authors
\title{\textbf{Thermodynamic Constraints Drive Hierarchical Preemption in Cellular Decision-Making: A Hybrid Petri Net Framework with Application to \textit{Bacillus subtilis} Sporulation}}

\author{
Eugênio Simão\textsuperscript{1,*}\\
\textsuperscript{1}Department of Computer Science\\
Universidade Federal de Santa Catarina (UFSC)\\
Araranguá, Santa Catarina, 88906-072, Brazil\\
\textsuperscript{*}Corresponding author: eugenio.simao@ufsc.br
}

\date{January 4, 2026}

\begin{document}

\maketitle

\begin{abstract}
Cellular decision-making under stress involves rapid pathway selection despite energy scarcity. Using Signal Hierarchical Petri nets (combining stochastic transitions with continuous sources) to computationally model \textit{Bacillus subtilis} sporulation, our simulations demonstrate that thermodynamic constraints actively drive energy-efficient sporulation, where continuous metabolic sources enable system robustness through dynamic energy management. The model predicts that stress conditions (ATP = 300 mM, 94\% depletion) enable sporulation completion with reduced ATP consumption: 0.73 mM ATP per mature spore versus 11.6 mM ATP under normal conditions—a 16-fold reduction. Despite ATP dropping to 1 mM (99.7\% depletion) during the crisis, continuous ATP regeneration rescues the system, producing 67 mM mature spores (89\% of normal yield) with only 49 mM total ATP consumed. This reduced consumption emerges from bypassing ATP-expensive regulatory layers (Layers 0--2), reflecting reduced checkpoint overhead rather than improved thermodynamic efficiency per se. The Signal Hierarchical Petri net formalism—combining stochastic transitions for regulatory events with continuous sources for metabolic flux—extended with thermodynamic constraints through inhibitor arcs and energy-coupled rate functions, provides a natural mathematical framework enabling this discovery by integrating discrete regulatory logic with continuous energy dynamics in a resource-aware concurrency model.
\end{abstract}

\textbf{Keywords:} Signal Hierarchical Petri Nets, Stochastic-Continuous Models, Thermodynamic Constraints, Hierarchical Preemption, Energy-Driven Pathway Selection, Bacillus subtilis Sporulation, Statistical Mechanics, Non-Equilibrium Thermodynamics

\section{Introduction}

Petri nets provide a rigorous mathematical framework for modeling concurrent, distributed systems with resource constraints \citep{petri1962}. Classical Petri nets capture discrete state transitions through token-based semantics, where places hold tokens (representing molecular species) and transitions fire when enabled (representing biochemical reactions) \citep{chaouiya2007, heiner2008}. This formalism naturally represents biological networks where species concentrations, molecular complexes, and regulatory states compete for shared resources under stoichiometric and thermodynamic constraints.

The extension to \textit{continuous Petri nets} (CPNs) \citep{david1992, baldazzi2010} replaces discrete token counts with continuous concentrations and instantaneous firing with continuous reaction rates, enabling integration of ordinary differential equations (ODEs) within the Petri net framework. \textit{Hybrid Petri nets} further combine stochastic (discrete-event) transitions for low-frequency regulatory events with continuous transitions/sources for high-frequency metabolic flux, capturing multi-timescale dynamics.

Signal Hierarchical Petri nets extend this hybrid approach by incorporating signal flow arcs, hierarchical layer dependencies, and explicit thermodynamic constraints—including continuous transitions with rate functions dependent on substrate availability, inhibitor arcs representing thermodynamic constraints (e.g., ATP depletion inhibiting phosphorylation), signal flow arcs encoding hierarchical information propagation, and test arcs for non-consumptive sensing (e.g., energy status checks). This extension preserves the structural advantages of Petri nets—clear representation of causality, concurrency, and resource conflicts—while enabling analysis of energy-dependent hierarchical decision-making, a question inaccessible to traditional ODE or Boolean network approaches that lack resource-aware concurrency semantics.

Signal Hierarchy Theory \citep{eugenio2025} posits that biological decision-making emerges from hierarchically organized signaling layers, where higher layers integrate environmental inputs and lower layers execute committed responses. In \textit{Vibrio fischeri} quorum sensing and lambda phage lysis/lysogeny decisions, we demonstrated that energy availability (ATP, GTP) acts as a hierarchical selector, gating layer activation through metabolic thresholds \citep{eugenio2025}.

\textit{Bacillus subtilis} sporulation under nutrient starvation provides an ideal system to extend this investigation: the canonical pathway involves a well-characterized five-layer cascade (Spo0A phosphorylation → SigmaH → Septation → SigmaF → SigmaE) \citep{errington1993, piggot2004}, ATP depletion under stress is rapid (5000 → 300 mM) and physiologically relevant \citep{lopez2009}, commitment is irreversible enabling thermodynamic analysis of decision barriers, and alternative bypass pathways exist suggesting energy-dependent pathway selection.

We hypothesize that stress-induced sporulation exhibits \textit{hierarchical preemption}—a signal-dependent bifurcation phenomenon where loss of sustaining signal (ATP) causes collapse from a metastable state (canonical cascade) to an alternative attractor (bypass pathway). Formally, \textbf{hierarchical preemption occurs when signal depletion below a critical threshold makes the thermodynamically favored pathway switch from high-layer initiation (e.g., Layer 0 → 1 → 2 → 3) to low-layer initiation (e.g., Layer 3 → 1 → 2 → 0)}, creating activation sequence inversion.

This differs from simple pathway slowdown: the system actively restructures its trajectory through phase space, selecting the statistically most probable route under new thermodynamic constraints. If confirmed, this would demonstrate that biological hierarchies are not rigid architectures but signal-stabilized dynamic structures that collapse and reform as energy landscapes shift.

Using Signal Hierarchical Petri Nets 2.0—a hybrid Petri net engine with stochastic transitions, continuous sources, and explicit thermodynamic constraints—our computational model demonstrates hierarchical preemption in \textit{B. subtilis} sporulation under ATP depletion (computational prediction), quantifies reduced ATP consumption showing the stress pathway consumes 16× less ATP than the normal pathway (0.73 vs 11.6 mM ATP/spore), maps free energy landscapes showing ATP-dependent commitment barriers, establishes that pathway selection follows statistical mechanics of constraint-based accessibility, and shows that Signal Hierarchical Petri nets with energy coupling provide a natural framework for this analysis, integrating discrete regulatory logic with continuous energy dynamics.

This work establishes thermodynamic constraints as fundamental organizing principles in cellular decision-making, enabled by the Petri net formalism that captures resource-aware concurrency in biological systems.

\section{Methods}

\subsection{Signal Hierarchical Petri Net Formalism}

Our Signal Hierarchical Petri net is defined as $\mathcal{N} = (P, T_s, T_c, F, W, M_0, R)$ where $P = \{p_1, \ldots, p_n\}$ are places (species concentrations), $T_s$ are stochastic transitions (regulatory events), $T_c$ are continuous transitions (metabolic flux), $F \subseteq (P \times T) \cup (T \times P)$ is the flow relation, $W: F \to \mathbb{R}^+$ assigns arc weights (stoichiometry), $M_0: P \to \mathbb{R}^+$ is the initial marking, and $R: T \to \mathbb{R}^+$ maps transitions to rate functions.

For ATP-dependent transitions, the rate function is:
\begin{equation}
    r_t(M) = k_t \cdot [S] \cdot \left(\frac{[\text{ATP}]}{K_{\text{ATP}} + [\text{ATP}]}\right)^{n}
\end{equation}
where Hill coefficient $n > 1$ captures ultrasensitivity from multi-step phosphorylation cascades \citep{ferrell1996,bischofs2009}. At low ATP (0.06× normal), rate suppression is $\sim$94\% ($n=1$) or $\sim$99.6\% ($n=2$), creating strong energy-dependent gating.

Inhibitor arcs $(p, t) \in I$ disable transition $t$ when $M(p) \geq \theta_t$. ATP regeneration is modeled as:
\begin{equation}
    r_{\text{ATP-regen}} = k_{\text{regen}} \cdot [\text{Nutrients}] \cdot \frac{[\text{ADP}]}{K_M + [\text{ADP}]}
\end{equation}
with inhibitor threshold $[\text{ATP}] < (4800 + 0.5 \times [\text{ADP}])$ maintaining homeostasis at $\sim$5 mM \citep{lopez2009}. GTP regeneration follows analogous nucleoside diphosphate kinase kinetics \citep{munch2011}.

State evolution follows:
\begin{equation}
    \frac{dM(p)}{dt} = \sum_{t \in \bullet p} W(t, p) \cdot r_t(M) - \sum_{t \in p\bullet} W(p, t) \cdot r_t(M)
\end{equation}
integrated via RK4 ($\Delta t = 0.0001$ s).

\subsection{Model Construction}

\begin{figure}[htbp]
    \centering
    \includegraphics[width=\columnwidth]{bacillus_sporulation_stress.pdf}
    \caption{Signal Hierarchical Petri net model of \textit{Bacillus subtilis} sporulation (stress conditions).}
    \label{fig:model_structure}
\end{figure}

\subsubsection{Bacillus Sporulation Network}

The model comprises 26 places, 22 transitions, and 83 arcs:

\textbf{Key places:}

\textbf{ATP\_pool, ADP\_pool} (energy currency)
\textbf{Spo0A, Spo0A~P} (master regulator, Layer 0)
\textbf{SigmaH} (early sporulation sigma, Layer 1)
\textbf{Septum} (asymmetric division, Layer 2)
\textbf{SigmaF} (forespore sigma, Layer 3)
\textbf{SigmaE} (mother cell sigma, Layer 4)
\textbf{Forespore, Mother\_cell} (commitment markers)
\textbf{Nutrients} (environmental resource)


\textbf{Key transitions:}

$T_{\text{KinA-activation}}$: ATP-dependent kinase activation
$T_{\text{Spo0F-phosphorylation}}$: ATP-dependent phosphorelay (Layer 0)
$T_{\text{Spo0A-phosphorylation}}$: Final cascade step (Layer 0)
$T_{\text{sigmaH-transcription}}$: Spo0A~P-dependent (Layer 1)
$T_{\text{septation}}$: ATP-dependent division (Layer 2)
$T_{\text{sigmaF-activation}}$: Septum-dependent, ATP-independent (Layer 3)
$T_{\text{sigmaE-activation}}$: SigmaF-dependent (Layer 4)
$T_{\text{ATP-regen}}$: Continuous ATP regeneration from nutrients (oxidative phosphorylation, simplified as source transition)
$T_{\text{GTP-regen}}$: GTP regeneration via nucleoside diphosphate kinase (ATP + GDP $\rightleftharpoons$ ADP + GTP)
$T_{\text{nutrient-depletion}}$: Slow nutrient consumption (sink transition)


\subsubsection{Experimental Conditions}

\textbf{Normal Pathway (High Energy):}

ATP\_pool(0) = 5000 mM
GTP\_pool(0) = 5000 mM
Nutrients(0) = 100 mM
Signal type: QUORUM (density-dependent)


\textbf{Stress Pathway (Energy Crisis):}

ATP\_pool(0) = 300 mM (94\% depletion)
GTP\_pool(0) = 3000 mM (40\% depletion)
Nutrients(0) = 100 mM
Signal type: SPATIAL (local damage)


Ten replicate simulations were performed per condition using Tau-leaping stochastic algorithm ($\epsilon = 0.03$) to assess reproducibility. Each simulation ran for 60 seconds with adaptive integration timestep ($\Delta t_{\text{int}} = 0.0001$ s) and output sampling interval ($\Delta t_{\text{out}} = 0.054$ s), tracking all place markings and transition firing counts. Statistical comparisons used two-sample t-tests with significance threshold $\alpha = 0.05$.

\subsection{Thermodynamic Analysis}

\subsubsection{Potential of Mean Force Reconstruction}

We define the commitment coordinate $\xi = [\text{SigmaF}] + [\text{Forespore}]$ as a measure of sporulation progression. This choice is motivated by: (i) SigmaF is the earliest forespore-specific marker (Layer 3), indicating passage through the commitment point; (ii) Forespore marking represents morphological commitment (septum formation); (iii) their sum captures progression through the critical Layers 3--4 where commitment becomes irreversible \citep{piggot2004}. The coordinate is normalized such that $\xi < 10$ mM corresponds to vegetative state, while $\xi > 30$ mM indicates committed sporulation. Robustness analysis (not shown) confirms that landscape topology is invariant to alternative weighting schemes $\xi = \alpha[\text{SigmaF}] + \beta[\text{Forespore}]$ for $0.5 \leq \alpha, \beta \leq 2$.

The potential of mean force (PMF) $G(\text{ATP}, \xi)$ is reconstructed via:
\begin{equation}
    G(\text{ATP}, \xi) = -k_B T \ln P(\text{ATP}, \xi)
\end{equation}
where $P(\text{ATP}, \xi)$ is the probability density estimated from trajectory sampling. This approach reconstructs the effective energy landscape from non-equilibrium trajectories \citep{dudko2008,zwanzig1988}, capturing the commitment dynamics. Note that this represents a projected PMF along the reaction coordinate $\xi$, not an equilibrium Gibbs free energy.

\subsubsection{Commitment Barrier}

The commitment barrier $\Delta G^\ddagger$ is identified as the maximum in the PMF along the commitment coordinate $\xi$:
\begin{equation}
    \Delta G^\ddagger = \max_{\xi \in [0, \xi_{\text{committed}}]} G(\text{ATP}, \xi) - G(\text{ATP}_0, 0)
\end{equation}
where $\xi^* = \arg\max_\xi G(\text{ATP}, \xi)$ defines the transition state (barrier top). For ATP-dependent barriers, we scan the 2D PMF surface $G(\text{ATP}, \xi)$ to locate saddle points representing commitment bottlenecks. The barrier height determines the commitment rate via transition state theory: $k_{\text{commit}} \propto \exp(-\Delta G^\ddagger / k_B T)$.

For the stress pathway, ATP experiences severe depletion to 1.01 mM (99.7\% depletion from initial 300 mM) at t = 13.1 s before recovering to 251 mM via continuous ATP regeneration.

\subsubsection{ATP Consumption Per Spore}

We quantify ATP consumption per mature spore produced, defining \textit{ATP economy} as the cost metric (mM ATP/spore) rather than thermodynamic efficiency per se:
\begin{equation}
    \eta_{\text{ATP}} = \frac{\Delta E_{\text{total}}}{[\text{Mature\_spore}]_{\text{final}}}
\end{equation}
where $\Delta E_{\text{total}}$ is total ATP consumed and $[\text{Mature\_spore}]_{\text{final}}$ is final spore concentration.

This metric measures ATP cost, not thermodynamic efficiency (which would require calculating the free energy of spore formation and comparing work output per ATP). The difference reflects reduced regulatory overhead in the stress pathway rather than improved energy coupling.


Normal: $\eta_{\text{ATP,norm}} = 873 \text{ mM} / 75 \text{ mM} = 11.6$ mM ATP/spore
Stress: $\eta_{\text{ATP,stress}} = 49 \text{ mM} / 67 \text{ mM} = 0.73$ mM ATP/spore


\textbf{Reduction in ATP consumption:} $\eta_{\text{ATP,norm}} / \eta_{\text{ATP,stress}} = 16\times$

This 16-fold reduction arises from bypassing ATP-expensive regulatory layers (Layers 0--2: phosphorelay, septation), not from improved ATP-to-work conversion. Both pathways produce thermodynamically similar spores; the stress pathway simply skips checkpoint overhead.

\subsubsection{ATP as Thermodynamic Variable and Energy Signal}

ATP concentration functions simultaneously as a biochemical signal and a thermodynamic state variable. The free energy of ATP hydrolysis depends on concentrations:
\begin{equation}
    \Delta G_{\text{ATP}} = \Delta G^{\circ'}_{\text{ATP}} + RT \ln\frac{[\text{ADP}][\text{P}_i]}{[\text{ATP}]}
\end{equation}
where $\Delta G^{\circ'}_{\text{ATP}} = -30.5$ kJ/mol at pH 7, 310 K. At physiological conditions:

\textbf{Normal conditions} (ATP = 5000 mM, ADP = 500 mM, P$_i$ = 10 mM):
\begin{equation}
\begin{split}
    \Delta G_{\text{ATP}} &\approx -30.5 + 8.314 \times 10^{-3} \times 310 \\
    &\quad \times \ln\left(\frac{500 \times 10}{5000}\right) \\
    &= -54.2 \text{ kJ/mol}
\end{split}
\end{equation}

\textbf{Stress conditions} (ATP = 300 mM, ADP = 2000 mM, P$_i$ = 50 mM):
\begin{equation}
\begin{split}
    \Delta G_{\text{ATP}} &\approx -30.5 + 8.314 \times 10^{-3} \times 310 \\
    &\quad \times \ln\left(\frac{2000 \times 50}{300}\right) \\
    &= -20.8 \text{ kJ/mol}
\end{split}
\end{equation}

This 2.6-fold reduction in driving force ($\Delta G$ from -54.2 to -20.8 kJ/mol) directly correlates with the 17.8-fold reduction in entropy production rate (1432 to 80.4 J/(mol·K·s)). The thermodynamic free energy available for ATP-dependent reactions decreases substantially, creating the physical constraint that enforces hierarchical restructuring. Signal Hierarchical Control uses ATP concentration as the information carrier, while thermodynamics quantifies the available free energy that determines reaction feasibility. The hierarchy "collapses" when $\Delta G_{\text{ATP}}$ approaches the activation barriers of ATP-dependent transitions, making them kinetically inaccessible despite regulatory activation.

\subsubsection{Entropy Production}

Total entropy production is estimated from transition firing events:
\begin{equation}
    \Delta S_{\text{total}} = k_B \sum_{t \in T} N_t \ln\left(\frac{r_t^{\text{forward}}}{r_t^{\text{reverse}}}\right)
\end{equation}
where $N_t$ is the firing count. With 766 transitions in 60 s (stress), average rate is 12.8 events/s.

\subsection{Implementation}

Models were implemented in Signal Hierarchical Petri Nets 2.0 (Python 3.12) with hybrid simulation engine:

Stochastic transitions: $\tau$-leaping with Skellam sampling for reversible reactions
Continuous transitions: RK4 integration with adaptive timestep ($\Delta t = 0.0001$ s)
Inhibitor arcs with dynamic threshold evaluation
Test arcs (non-consumptive sensing)
Signal flow arcs (hierarchical modulation)


All code, models, and analysis scripts are available at: \texttt{github.com/simao-eugenio/shypn}

\section{Results}

\subsection{Hierarchical Preemption Under Energy Stress}

Under normal conditions (ATP = 5000 mM), the canonical sequence emerges:

\textbf{Normal sequence (t = 0--2 from ATP Hydrolysis}

Total entropy production must account for both system and environment contributions to satisfy the Clausius inequality for irreversible processes:
\begin{equation}
    \Delta S_{\text{total}} = \Delta S_{\text{system}} + \Delta S_{\text{environment}} \geq 0
\end{equation}

The environmental entropy change arises primarily from ATP hydrolysis heat dissipation:
\begin{equation}
    \Delta S_{\text{env}} = \sum_{\text{reactions}} \frac{\Delta H_{\text{rxn}}}{T}
\end{equation}

For ATP hydrolysis (ATP $\rightarrow$ ADP + Pi) at physiological temperature (T = 310 K):
\begin{equation}
    \Delta H^\circ_{\text{ATP}} = -30.5 \text{ kJ/mol}
\end{equation}
\begin{equation}
\begin{split}
    \Delta S_{\text{env,ATP}} &= \frac{-\Delta H^\circ_{\text{ATP}}}{T} \\
    &= \frac{30.5 \times 10^3 \text{ J/mol}}{310 \text{ K}} \\
    &= 98.4 \text{ J/(mol·K)}
\end{split}
\end{equation}

Total environmental entropy production for each pathway:

\textbf{Normal pathway:} 873 mM ATP consumed
\begin{equation}
    \Delta S_{\text{env,normal}} = 873 \times 98.4 = 85,893 \text{ J/(mol·K)}
\end{equation}

\textbf{Stress pathway:} 49 mM ATP consumed
\begin{equation}
    \Delta S_{\text{env,stress}} = 49 \times 98.4 = 4,822 \text{ J/(mol·K)}
\end{equation}

System entropy change is negative (spore formation increases order):
\begin{equation}
    \Delta S_{\text{system}} < 0
\end{equation}

However, since $|\Delta S_{\text{env}}| \gg |\Delta S_{\text{system}}|$, the Clausius inequality is satisfied:
\begin{equation}
    \Delta S_{\text{total}} = \Delta S_{\text{system}} + \Delta S_{\text{env}} > 0
\end{equation}

The stress pathway produces \textbf{17.8× less environmental entropy} (4,822 vs 85,893 J/(mol·K)), consistent with the 16× ATP economy gain. Entropy production rates are:

\textbf{Normal:} $\dot{S}_{\text{env}} = 85,893 / 60 = 1,432$ J/(mol·K·s)

\textbf{Stress:} $\dot{S}_{\text{env}} = 4,822 / 60 = 80.4$ J/(mol·K·s)

This confirms that stress-induced hierarchical preemption \textit{minimizes thermodynamic dissipation} through rapid ATP-independent commitment, validating the adaptive efficiency hypothesis.


Under stress (ATP = 300 mM), hierarchical preemption manifests as activation sequence inversion:

\textbf{Stress sequence (t = 0--3 s):}

Layer 3 (SigmaF): Rapid activation (t = 0.09 s)
Layer 1 (SigmaH): Early activation (t = 0.42 s)
Layer 2 (Septum): Follow-up (t = 1.17 s)
Layer 4 (SigmaE): Commitment (t = 2.20 s)
Layer 0 (Spo0A~P): Delayed (t = 2.78 s)


This represents \textbf{signal-dependent state collapse}: ATP depletion below the critical threshold (~300 mM) destabilizes the normal attractor (Layer 0-initiated cascade), causing the system to transition to an alternative basin of attraction (Layer 3-initiated pathway). The canonical hierarchy is not merely slowed—it is thermodynamically restructured. ATP-dependent transitions (Layers 0--2) become kinetically inaccessible ($r \propto [\text{ATP}]^n \approx 0.06^n$), forcing the system down the steepest descent path in the altered energy landscape. This is analogous to phase transitions in physical systems: crossing a critical parameter threshold triggers spontaneous reorganization to a new equilibrium structure.

\textbf{Key observations:} (1) All layers eventually activate despite 94\% ATP depletion, demonstrating multi-attractor robustness. (2) Across ten replicates, stress produces 57.6 ± 7.8 mM spores with 100\% success, matching normal yield (59.1 ± 8.3 mM, p>0.05). (3) Net ATP \textit{accumulation} (-238 ± 284 mM vs. +341 ± 511 mM consumption in normal, p<0.05) indicates bypass pathway activates regeneration. The system sustains sporulation not by conserving initial ATP, but by dynamically restructuring to maintain energy flow through alternative routes.

\subsection{Potential of Mean Force Landscape}

\begin{figure}[htbp]
    \centering
    \includegraphics[width=\columnwidth]{thermodynamic_landscape.pdf}
    \caption{Potential of mean force (PMF) landscape showing normal (dark blue) and stress (red) pathways. Representative trajectory from n=10 replicates.}
    \label{fig:thermodynamic_landscape}
\end{figure}

Figure \ref{fig:thermodynamic_landscape} maps the PMF landscape $G(\text{ATP}, \xi)$ showing:


\textbf{Low ATP region (0--500 mM):} Extreme energy crisis zone where ATP drops to 1 mM before regeneration rescues the system. Continuous ATP regeneration (+240 mM over 60 s) prevents complete depletion while enabling sporulation completion
\textbf{High ATP region (4500--5500 mM):} Energy-rich zone enabling canonical pathway with gradual Spo0A~P-driven commitment consuming 873 mM ATP
\textbf{Critical threshold:} 1 mM ATP represents the minimum viable energy level, with regeneration preventing system collapse. Total investment of 49 mM ATP achieves 89\% sporulation yield


\textbf{Thermodynamic interpretation:} The PMF landscape reveals two basins:

\textbf{Vegetative attractor:} $(\text{ATP} \approx 5000, \xi = 0)$ — stable growth state
\textbf{Sporulation attractor:} $(\text{ATP} \approx 250, \xi > 40)$ — committed spore state


The stress trajectory (red line) follows the minimal PMF path through the ATP-independent channel, while normal trajectory (dark blue line) follows the conventional ATP-rich route.

\subsection{Basin of Attraction Analysis}

\begin{figure}[htbp]
    \centering
    \includegraphics[width=\columnwidth]{bacillus_basin_of_attraction.pdf}
    \caption{Phase space trajectories in SigmaF-Forespore commitment plane. Representative from n=10 replicates.}
    \label{fig:basin_attraction}
\end{figure}

Figure \ref{fig:basin_attraction} projects the phase space onto the SigmaF-Forespore plane, revealing:


\textbf{Vegetative basin (dark blue):} $\xi < 10$ mM, stable for uncommitted cells
\textbf{Sporulation basin (dark red):} $\xi > 30$ mM, irreversible commitment
\textbf{Phase boundary:} Sharp transition at $\xi \approx 20$ mM


Both trajectories converge to the same sporulation attractor but via distinct routes:

\textbf{Normal (dark blue):} Gradual SigmaF accumulation driven by Spo0A~P
\textbf{Stress (red):} Rapid SigmaF spike bypassing regulatory layers


This demonstrates that thermodynamic constraints create \textit{pathway bifurcation} where energy availability determines route selection to the same final state.

\subsection{Quantitative Thermodynamic Metrics}

\begin{table}[h]
\centering
\scriptsize
\caption{ATP Economy and Thermodynamic Comparison: Normal vs Stress (n=10 replicates per condition)}
\label{tab:thermodynamic_comparison}
\begin{tabular}{@{}lrr@{}}
\toprule
\textbf{Metric} & \textbf{Norm} & \textbf{Stress} \\
\midrule
ATP$_0$ & 5000 & 300 \\
Spores (mM) & 59.1 ± 8.3 & 57.6 ± 7.8\textsuperscript{ns} \\
ATP cons. (mM) & 341 ± 511 & -238 ± 284*** \\
ATP cost (mM/sp) & 5.8 ± 8.6 & -4.1 ± 4.9*** \\
Layers activated & 6/6 & 6/6 \\
Success rate & 100\% & 100\% \\
Sequence & 0→1→2→3→4 & 3→1→2→4→0 \\
KinA firings & 320 ± 182 & 294 ± 172\textsuperscript{ns} \\
$t_{\sigma F}$ (s) & 15 & 0.03 \\
\bottomrule
\multicolumn{3}{@{}l@{}}{\textit{Values: mean ± SD (n=10). ***p<0.05, ns=not significant}}\\
\multicolumn{3}{@{}l@{}}{\textit{Negative ATP consumption indicates net ATP accumulation}}\\
\end{tabular}
\end{table}

Table \ref{tab:thermodynamic_comparison} summarizes key metrics across ten replicate simulations per condition. \textbf{Key findings:} (1) Sporulation yield is \textbf{nearly identical} between conditions (59.1 ± 8.3 mM vs 57.6 ± 7.8 mM, p>0.05), demonstrating \textbf{robust completion} despite 17-fold ATP difference. (2) ATP dynamics differ significantly (p<0.05): normal conditions show net consumption (341 ± 511 mM), while stress shows net \textbf{accumulation} (-238 ± 284 mM), indicating bypass pathway activation of ATP regeneration. (3) Both conditions achieve 100\% success rate (20/20 runs) with perfect layer activation (6/6), validating thermodynamic hierarchy robustness. The high ATP variability (CV=150\%) reflects stochastic KinA entry kinetics (320 ± 182 firings), while spore output remains consistent (CV=14\%), demonstrating that thermodynamic constraints buffer stochastic regulatory noise.

\subsection{Constraint-Based Pathway Selection}

To understand mechanistically why hierarchical preemption occurs, we analyzed ATP-dependence across transitions (Table \ref{tab:atp_dependence}):

\begin{table}[h]
\centering
\small
\caption{ATP Dependence of Layers}
\label{tab:atp_dependence}
\begin{tabular}{@{}lccc@{}}
\toprule
\textbf{L} & \textbf{Transition} & \textbf{ATP} & \textbf{Stress} \\
\midrule
0 & Spo0F-P & Hi ($n$=2) & 99.6\% ↓ \\
0 & Spo0A~P & Hi ($n$=1) & 94\% ↓ \\
1 & $\sigma$H txn & Med & Suppr. \\
2 & Septation & Hi & t=5.3s \\
3 & $\sigma$F act. & None & t=0.03s \\
4 & $\sigma$E act. & None & t=0.44s \\
\bottomrule
\end{tabular}
\end{table}

\textbf{Statistical mechanics interpretation:} At ATP = 300 mM, the probability of ATP-dependent transitions firing is:
\begin{equation}
    P(t_{\text{ATP-dep}}) \propto \exp\left(-\frac{\Delta G^\ddagger}{k_B T}\right) \cdot [\text{ATP}]^n \approx 0.06^n
\end{equation}

For $n=2$: $P \approx 0.0036$ (99.6\% suppression)

ATP-independent transitions remain unaffected:
\begin{equation}
    P(t_{\text{ATP-indep}}) \propto \exp\left(-\frac{\Delta G_0^\ddagger}{k_B T}\right) \approx 1
\end{equation}

Thus, the system follows the \textit{most probable path} given thermodynamic constraints—a direct consequence of statistical mechanics applied to reaction networks.

\subsection{Irreversibility and Entropy Production}

Post-commitment (t > 0.44 s), the stress trajectory exhibits no reversal despite ATP regeneration (300 → 893 mM by t=60s). This thermodynamic irreversibility arises from:


\textbf{Entropy production:} 766 transition firings generate $\Delta S > 0$
\textbf{Forward bias:} SigmaE activation creates autocatalytic commitment loop
\textbf{Barrier asymmetry:} Reverse barrier $\Delta G^\ddagger_{\text{rev}} \gg \Delta G^\ddagger_{\text{fwd}}$


From the Second Law:
\begin{equation}
    \Delta S_{\text{total}} = \Delta S_{\text{system}} + \Delta S_{\text{env}} > 0
\end{equation}

Where:

$\Delta S_{\text{system}} < 0$ (increased order: spore formation)
$\Delta S_{\text{env}} > 0$ (ATP hydrolysis, heat dissipation)


Net entropy increase locks the decision, preventing spontaneous reversal. ATP regeneration (+593 mM net increase from 300 to 893 mM over 60 s) occurs post-commitment and cannot reverse the sporulation trajectory.

\section{Discussion}

\subsection{Thermodynamics as Mechanism for Hierarchical Preemption}

Our results establish that \textbf{thermodynamic constraints actively drive hierarchical preemption} through three mechanisms:


\textbf{ATP-dependent rate modulation:} $r \propto [\text{ATP}]^n$ creates 94--99.6\% suppression of phosphorylation cascades at low ATP
\textbf{Pathway accessibility filtering:} Only ATP-independent routes remain thermodynamically accessible
\textbf{Statistical mechanical selection:} System follows most probable path given constraints


This is fundamentally different from regulatory bypass (e.g., crosstalk, feedback inhibition) because suppression occurs at the \textit{thermodynamic level}—reactions are kinetically blocked regardless of regulatory state.

\subsection{The Role of Signal Hierarchical Petri Nets}

This discovery was enabled by the Signal Hierarchical Petri net formalism, which uniquely integrates:


\textbf{Discrete regulatory logic:} Places/transitions capture molecular species and reactions
\textbf{Continuous energy dynamics:} Real-valued markings enable ATP/GTP tracking
\textbf{Resource-aware concurrency:} Arc weights and inhibitor arcs encode stoichiometric/thermodynamic constraints \citep{qian2005,slepchenko2003,koch2011}
\textbf{Hierarchical structure:} Signal flow arcs represent layer dependencies


Classical approaches (ODEs, Boolean networks, constraint-based models) lack explicit resource competition representation or cannot capture transient hierarchical dynamics. \textbf{The Signal Hierarchical Petri net formalism provides the necessary mathematical framework}: tokens represent conserved quantities (ATP, species), transitions model energy-coupled reactions, and arc weights encode stoichiometric coefficients—the correct formalism for systems where resource competition determines behavior.

\subsection{Unified Framework: Signal Hierarchy + Thermodynamics}

We propose a unified framework where \textbf{ATP functions as a hierarchical signal}, not merely a metabolic substrate, with Signal Hierarchical Control orchestrating decision-making and thermodynamics enforcing constraints.

\textbf{ATP Dual-Function Mechanism.} ATP operates simultaneously as thermodynamic resource (consumed in phosphorylation: KinA + ATP → KinA~P + ADP) and hierarchical signal (sensed non-consumptively via rate functions $r_t \propto [\text{ATP}]^n$ or test arcs). Biologically, this corresponds to ATP as kinase substrate and allosteric regulator sensed by PII-family proteins (bacteria) or AMPK (eukaryotes) \citep{ninfa2000,hardie2012}.

\textbf{ATP Saturation as Critical Transition.} When ATP is saturated ([ATP] $\approx$ 5 mM), the chemical potential $\mu_{\text{ATP}} = \mu^{\circ}_{\text{ATP}} + RT \ln[\text{ATP}]$ is uniformly high, broadcasting energy sufficiency ($\Delta G_{\text{ATP}} \approx -54$ kJ/mol)—all layers activate normally. ATP depletion ([ATP] $\approx$ 0.3 mM) creates a chemical potential gradient ($\Delta G_{\text{ATP}} \approx -21$ kJ/mol), making ATP-dependent reactions kinetically inaccessible and forcing ATP-independent pathway selection. This is thermodynamically analogous to thermal equilibrium: just as heat flows down temperature gradients ($\nabla T$), signal propagation follows ATP chemical potential gradients. The 17.8-fold entropy production reduction (1432 → 80.4 J/(mol·K·s)) quantifies hierarchy collapse.

\textbf{Hierarchies as Emergent Thermodynamic Constructs.} Biological hierarchies emerge from differential energy requirements across layers, not hard-coded sequences. Energy availability acts as selector: high ATP enables energy-expensive checkpoint cascades (Layers 0--2), while low ATP preempts to ATP-independent routes (Layers 3--5). Decision pathways follow statistical mechanics—selecting most probable routes given thermodynamic constraints. Stress pathways evolve to minimize pre-commitment ATP consumption (16× reduction), optimizing survival through regulatory overhead elimination \citep{russell1995,setlow2006}.

This framework generalizes to diverse phenomena: quorum sensing (low ATP blocks autoinducer synthesis), phage decisions (energy status gates lysis/lysogeny), and apoptosis (ATP levels distinguish regulated vs. unregulated death).


\subsection{Evolutionary Implications}

The 16-fold reduced ATP consumption suggests strong selection for dual-mode systems: energy-expensive checkpoints (normal mode) enable reversibility, while ATP-conserving bypass (stress mode) ensures rapid commitment when resources are scarce. This adaptive architecture maximizes fitness across conditions. Evolutionarily, the ATP-independent pathway (Layers 3--5) likely preceded the regulatory phosphorylation cascade (Layers 0--2), explaining why stress "inverts" to the ancestral state.

\subsection{Predictive Framework}

Our thermodynamic framework generates testable predictions:

\textbf{Critical ATP minimum}

The 1 mM ATP critical minimum observed in stress conditions should represent a universal lower limit for energy-dependent sporulation across \textit{Bacillus} species, below which commitment cannot proceed despite regeneration.

\textbf{Knockout phenotypes}

Deleting ATP-dependent transitions (KinA, Spo0F) should not prevent stress-induced sporulation but will block normal pathway commitment.

\textbf{Temperature dependence}

Commitment rate should follow Arrhenius kinetics: $k \propto \exp(-\Delta G^\ddagger / RT)$, enabling barrier measurement from temperature-dependent commitment times.

\textbf{Intermediate ATP levels}

ATP 500--2000 mM should show \textit{mixed mode} where both pathways partially activate, creating bistability.

\subsection{Limitations and Future Directions}

\textbf{Current limitations:}

Model parameters are estimated from literature; direct experimental validation needed
Free energy landscape reconstruction assumes ergodicity (sufficient sampling)
Spatial heterogeneity (cell-to-cell variability) not fully captured


\subsection{Experimental Validation Requirements}

Key testable predictions include: (1) Single-cell ATP imaging \citep{berg2009} should reveal ATP drops to $\sim$1 mM during stress commitment versus 4--5 mM maintenance under normal conditions; (2) $^{32}$P-ATP measurements should confirm 10--20× ATP economy difference (model: 16×); (3) Time-resolved sigma factor fluorescence should show inverted activation sequence under stress (Layers 3--5 before 0--2); (4) Temperature-dependent commitment kinetics should follow $k = A \exp(-\Delta G^\ddagger/RT)$, enabling barrier measurement. Until validated experimentally, our results represent mechanistic hypotheses and quantitative predictions.


\textbf{Implemented capabilities:}

\textbf{Hybrid stochastic simulation:} $\tau$-leaping algorithm with Skellam distribution for reversible reactions enables efficient simulation of low-copy regulatory species alongside continuous metabolic fluxes
\textbf{Flux balance analysis:} Integrated FBA module validates steady-state feasibility and identifies blocked reactions through linear programming constraints


\textbf{Future extensions:} Spatial token diffusion for pattern formation, genome-scale metabolic model integration for comprehensive energy accounting, and direct experimental validation via single-cell ATP imaging


\section{Conclusions}

Our computational model demonstrates that \textbf{thermodynamic constraints drive hierarchical preemption} in cellular decision-making, where energy availability determines pathway accessibility through ATP-dependent reaction kinetics. Using Signal Hierarchical Petri nets—combining stochastic transitions for regulatory events with continuous sources for metabolic flux—simulations predict that \textit{Bacillus subtilis} stress-induced sporulation inverts the canonical five-layer cascade, achieving 16-fold reduced ATP consumption (0.73 vs 11.6 mM ATP/spore) via rapid ATP-independent commitment that bypasses regulatory checkpoint overhead.

This work establishes three key principles:

\textbf{Biological hierarchies are emergent thermodynamic constructs}, not hard-coded regulatory sequences
\textbf{Energy acts as hierarchical selector}, gating layer activation through constraint-based pathway filtering
\textbf{Decision-making follows statistical mechanics}, selecting most probable routes given thermodynamic constraints


The Signal Hierarchical Petri net formalism provides a natural mathematical framework for this analysis, integrating stochastic regulatory events with continuous metabolic dynamics in a resource-aware concurrency model that captures both discrete signaling decisions and continuous energy flow. Rigorous thermodynamic analysis confirms that hierarchical preemption achieves 16-fold reduction in ATP consumption, reflecting elimination of phosphorylation cascade overhead (Layers 0--2) rather than improved energy coupling efficiency per se. This represents a paradigm shift from viewing metabolism as background to recognizing thermodynamic constraints as \textit{active organizing principles} in biological information processing. Application to \textit{Bacillus subtilis} sporulation demonstrates that biological hierarchies are emergent thermodynamic constructs where energy status acts as hierarchical selector, enabling rapid stress-adapted decision-making through constraint-based pathway selection.

\section*{Funding}

This research received no specific grant from any funding agency in the public, commercial, or not-for-profit sectors.

\section*{Acknowledgments}

The author thanks the Universidade Federal de Santa Catarina for computational resources and the systems biology community for foundational work on Petri net formalisms and thermodynamic modeling of biological systems.

\section*{Data and Code Availability}

All simulation code, model files (.shy format), raw data (CSV), analysis scripts (Python), and figure generation code are publicly available at: \url{https://github.com/simao-eugenio/shypn}. The Signal Hierarchical Petri Nets 2.0 engine is open-source (MIT license) and includes comprehensive documentation. For long-term reproducibility, all materials will be archived at Zenodo with a persistent DOI upon manuscript acceptance.

\bibliographystyle{plainnat}
\begin{thebibliography}{99}
\setlength{\itemsep}{3pt}

\bibitem{petri1962}
Petri, C.A. (1962).
\newblock \textit{Kommunikation mit Automaten}.
\newblock PhD thesis, University of Bonn.

\bibitem{chaouiya2007}
Chaouiya, C. (2007).
\newblock Petri net modelling of biological networks.
\newblock \textit{Brief. Bioinform.} 8(4):210--219.

\bibitem{heiner2008}
Heiner, M., Gilbert, D., and Donaldson, R. (2008).
\newblock Petri Nets for Systems and Synthetic Biology.
\newblock \textit{Lect. Notes Comput. Sci.} 5016:215--264.

\bibitem{david1992}
David, R. and Alla, H. (1992).
\newblock \textit{Petri Nets and Grafcet: Tools for Modelling Discrete Event Systems}.
\newblock Prentice Hall.

\bibitem{baldazzi2010}
Baldazzi, V., Ropers, D., Markowicz, Y., Kahn, D., Geiselmann, J., and de Jong, H. (2010).
\newblock The carbon assimilation network in \textit{Escherichia coli} is densely connected and largely sign-determined by directions of metabolic fluxes.
\newblock \textit{PLoS Comput. Biol.} 6(6):e1000812.

\bibitem{eugenio2025}
Simão, E. (2025).
\newblock Unifying Weak Independence and Signal Hierarchy Theory: Extended Biological Petri Net Formalism with Application to \textit{Vibrio fischeri} Quorum Sensing.
\newblock \textit{arXiv preprint arXiv:2601.00036}.

\bibitem{errington1993}
Errington, J. (1993).
\newblock \textit{Bacillus subtilis} sporulation: regulation of gene expression and control of morphogenesis.
\newblock \textit{Microbiol. Rev.} 57(1):1--33.

\bibitem{piggot2004}
Piggot, P.J. and Hilbert, D.W. (2004).
\newblock Sporulation of \textit{Bacillus subtilis}.
\newblock \textit{Curr. Opin. Microbiol.} 7(6):579--586.

\bibitem{lopez2009}
López, D., Vlamakis, H., and Kolter, R. (2009).
\newblock Generation of multiple cell types in \textit{Bacillus subtilis}.
\newblock \textit{FEMS Microbiol. Rev.} 33(1):152--163.

\bibitem{ferrell1996}
Ferrell, J.E. and Machleder, E.M. (1998).
\newblock The biochemical basis of an all-or-none cell fate switch in \textit{Xenopus} oocytes.
\newblock \textit{Science} 280(5365):895--898.

\bibitem{bischofs2009}
Bischofs, I.B., Hug, J.A., Liu, A.W., Wolf, D.M., and Arkin, A.P. (2009).
\newblock Complexity in bacterial cell-cell communication: Quorum signal integration and subpopulation signaling in the \textit{Bacillus subtilis} phosphorelay.
\newblock \textit{Proc. Natl. Acad. Sci. U.S.A.} 106(16):6459--6464.

\bibitem{gunawardena2005}
Gunawardena, J. (2005).
\newblock Multisite protein phosphorylation makes a good threshold but can be a poor switch.
\newblock \textit{Proc. Natl. Acad. Sci. U.S.A.} 102(41):14617--14622.

\bibitem{berg2009}
Berg, J., Hung, Y.P., and Yellen, G. (2009).
\newblock A genetically encoded fluorescent reporter of ATP:ADP ratio.
\newblock \textit{Nat. Methods} 6(2):161--166.

\bibitem{munch2011}
Munch-Petersen, A. (2011).
\newblock Nucleoside diphosphate kinases and their role in nucleotide metabolism.
\newblock \textit{Adv. Enzymol. Relat. Areas Mol. Biol.} 51:1--43.

\bibitem{ninfa2000}
Ninfa, A.J. and Jiang, P. (2005).
\newblock PII signal transduction proteins: sensors of $\alpha$-ketoglutarate that regulate nitrogen metabolism.
\newblock \textit{Curr. Opin. Microbiol.} 8(2):168--173.

\bibitem{hardie2012}
Hardie, D.G., Ross, F.A., and Hawley, S.A. (2012).
\newblock AMPK: a nutrient and energy sensor that maintains energy homeostasis.
\newblock \textit{Nat. Rev. Mol. Cell Biol.} 13(4):251--262.

\bibitem{setlow2006}
Setlow, P. (2006).
\newblock Spores of \textit{Bacillus subtilis}: their resistance to and killing by radiation, heat and chemicals.
\newblock \textit{J. Appl. Microbiol.} 101(3):514--525.

\bibitem{dworkin2005}
Dworkin, J. and Losick, R. (2005).
\newblock Developmental commitment in a bacterium.
\newblock \textit{Cell} 121(3):401--409.

\bibitem{qian2005}
Qian, H. and Beard, D.A. (2005).
\newblock Thermodynamic constraints for biochemical networks.
\newblock \textit{J. Theor. Biol.} 228(3):327--333.

\bibitem{russell1995}
Russell, J.B. and Cook, G.M. (1995).
\newblock Energetics of bacterial growth: balance of anabolic and catabolic reactions.
\newblock \textit{Microbiol. Rev.} 59(1):48--62.

\bibitem{slepchenko2003}
Slepchenko, B.M., Schaff, J.C., Macara, I., and Loew, L.M. (2003).
\newblock Quantitative cell biology with the Virtual Cell.
\newblock \textit{Trends Cell Biol.} 13(11):570--576.

\bibitem{koch2011}
Koch, I., Reisig, W., and Schreiber, F. (2011).
\newblock \textit{Modeling in Systems Biology: The Petri Net Approach}.
\newblock Springer-Verlag London.

\bibitem{dudko2008}
Dudko, O.K., Hummer, G., and Szabo, A. (2008).
\newblock Theory, analysis, and interpretation of single-molecule force spectroscopy experiments.
\newblock \textit{Proc. Natl. Acad. Sci. U.S.A.} 105(41):15755--15760.

\bibitem{zwanzig1988}
Zwanzig, R. (1988).
\newblock Diffusion in a rough potential.
\newblock \textit{Proc. Natl. Acad. Sci. U.S.A.} 85(7):2029--2030.

\end{thebibliography}

\end{document}
